
\subsubsection{6.1 Key Findings Interpretation}

The H-C1 feasibility scan reveals that doctest-passing filtering is structurally infeasible
at 500M-token scale in clean subprocess environments. The $31\times$ gap has a clear mechanistic
explanation: import isolation. Python library code is designed for installed environments,
not clean subprocess execution. The \texttt{import numpy as np} at the top of a doctest is an
implicit dependency on a pre-installed package, not a self-contained executable example.
Import errors constitute the large majority of Phase C failures, consistent with this
characterization.

The Phase A-to-B attrition (3.1\% to 2.0\%, approximately 1.$5\times )$ reflects two factors:
files containing syntax errors that prevent AST parsing, and \texttt{>>>} occurrences outside
docstrings (e.g., shell prompts in comments). The Phase B-to-C attrition (2.0\% to 0.1\%,
approximately $20\times )$ is substantially larger and is consistent with the import isolation
hypothesis.

Compile-only filtering is feasible at scale and avoids the import isolation problem that
makes doctest-passing filtering impractical for raw Python SFT corpus construction. Compile-only
filtering applies a syntax validity gate, retains a substantial fraction of corpus files, and
requires no dependency management. Whether this syntax-validity gate yields measurable SFT
quality improvement over existing heuristic filters remains to be confirmed by H-E1.

\subsubsection{6.2 Limitations}

\textbf{L1: Dataset Fallback.} We used \texttt{codeparrot/codeparrot-clean-valid} as a proxy for
\texttt{bigcode/the-stack-dedup} (access-gated; confirmed in \texttt{experiment.log}). The $31\times$ gap is
robust to reasonable dataset composition differences: even at $5\times$ higher rate (0.5\%), the
token pool (0.02M) remains 25,$000\times$ below the 500M target. The PIVOT conclusion holds under
any plausible dataset composition variation.

\textbf{L2: H-E1 Pending.} The SFT training experiment has not been executed. The compile-only SFT
benefit claim is supported by empirical precedent (phi-1, EffiCoder) but requires confirmation
via H-E1.

\textbf{L3: Mechanism Disambiguation.} Planned mechanism experiments (perplexity comparison,
n-gram overlap analysis) are not executed. The mechanism by which compile-only filtering
could improve benchmark performance remains hypothesized.

\textbf{L4: Sample Size.} The scan covers 10,000 files, consistent with The Stack paper's own
analysis methodology \citep{Kocetkov2022}. The estimate of ~12,960 executable-doctest
files in the full corpus extrapolates from this sample; the confidence interval on the 0.1\%
rate at n=10,000 (binomial 95\% CI: approximately 0.05\%--0.18\%) does not alter the PIVOT
conclusion.

\subsubsection{6.3 Broader Impact}

This work is directly relevant to teams building open-source code LLM training pipelines.
The validated Phase A/B/C pipeline is reusable for corpus characterization studies. The
SFT quality benefit of compile-only filtering will be confirmed by H-E1.

\subsubsection{6.4 Future Work}

\textbf{(1)} Execute H-E1: Qwen2.5-Coder-1.5B, compile-only vs. unfiltered, 500M tokens,
HumanEval/MBPP pass@1.
\textbf{(2)} Re-run Phase C with pre-installed scientific Python packages (numpy, scipy, pandas,
torch) to quantify the contribution of import isolation vs. stale documentation to Phase C
failures.
\textbf{(3)} Re-run the scan on \texttt{bigcode/the-stack-dedup} once access is obtained to directly
characterize the intended target corpus.
\textbf{(4)} Multi-language generalization of compile-only filtering (Rust, Go, TypeScript).

